\subsection*{Synchronous and Asynchronous}

\textbf{Asynchronous DMA}, or \textbf{Async-DMA}, is also known as \textbf{Non-Blocking DMA} because it does not block the execution of subsequent code. Conversely, \textbf{Synchronous DMA}, or \textbf{Sync-DMA}, is referred to as \textbf{Blocking DMA} because the code does not proceed until the DMA operation is complete. These terms will be used interchangeably in the following text.

In Croqtile, the \texttt{.async} suffix in a \textbf{DMA operation} indicates that the operation is asynchronous. The following code provides an example:

\begin{lstlisting}[language=croqtile,caption={Synchronous and Asynchronous}]
f0 = dma.copy.async data0 => shared; // non-blocking
f1 = dma.copy data1 => shared;       // blocking
...
wait f0;                             // synchronization
\end{lstlisting}

In this example, the DMA operation that produces the future \texttt{f0} is asynchronous, allowing the program to continue executing while the data transfer occurs in the background. The DMA operation that produces \texttt{f1} is synchronous, meaning the program will pause until the data is fully transferred from \texttt{data1} to \texttt{shared} storage. Meanwhile, the data copy of \texttt{data0} might still be in progress.

Executing the statement \texttt{wait f0;} causes the program to check if \texttt{data0} has been fully transferred to \texttt{shared} storage. This synchronizes the main execution flow with the asynchronous DMA of \texttt{f0}. Therefore programmers must make sure there is no "dangling future* left in your code.
